% !TeX root = ../../Book3.tex
\section{曲线与同伦}
\label{b3:sec:0901}

解析延拓沿路径传递局部函数。两条路径是否可以连续变形、变形时是否经过奇点，将决定它们传递到终点的函数芽是否相同。本节把这些拓扑数据固定下来。

\subsection{路径}
\label{b3:subsec:090101}

路径指连续映射 \(\gamma:[0,1]\rightarrow U\)。起终点为 \(\gamma(0),\gamma(1)\)，反向路径记 \(\gamma^{-1}(t)=\gamma(1-t)\)。若 \(\gamma\) 的终点是 \(\eta\) 的起点，拼接 \(\gamma*\eta\) 先走 \(\gamma\) 再走 \(\eta\)，分别在两半区间线性重参数化。

\ProseParagraphBreak
沿连续路径作解析延拓不需要曲线可求长，因为传递的是重叠邻域上的函数芽，而不是沿曲线积分。需要积分时，平面区域中的连续路径可以在固定端点下用有限折线路径逼近；把紧路径像覆盖于有限小圆盘中，逐段以线段替换便给出这种逼近及圆盘内同伦。

\subsection{同伦}
\label{b3:subsec:090102}

两个同起终点的路径端点固定同伦，指存在连续 \(H:[0,1]^2\rightarrow U\)，使 \(H(0,t)=\gamma(t)\)、\(H(1,t)=\eta(t)\)，且 \(H(s,0),H(s,1)\) 分别是固定端点。反射参数、拼接参数与恒定变形证明同伦是等价关系。

\ProseParagraphBreak
自由同伦允许闭路基点移动；它与端点固定同伦不同。基本群使用后一种。沿着移动基点的路径传递闭路类时，会出现共轭，而不是原样保留每个群元素。对整数绕数这项差别被交换性遮住，对两个以上孔洞则不能忽略。

\begin{proposition}\label{b3:prop:09-path-homotopies}
拼接在同伦类上满足结合律，常路径为单位，\(\gamma*\gamma^{-1}\) 同伦于常路径。
\end{proposition}
\begin{proof}
三个路径的不同括号仅改变各段占用的参数区间，连续移动分割点即给结合律同伦。常路径占用的参数区间可连续缩到长度0，给单位关系。往返路径可先只走 \(\gamma\) 到参数 \(1-s\)，再沿原路返回；当 \(s\) 从0变到1时，往返部分连续缩成起点，给逆元关系。
\end{proof}

\subsection{基本群}
\label{b3:subsec:090103}

\begin{definition}\label{b3:def:09-fundamental-group}
以 \(a\in U\) 为基点的闭路径端点固定同伦类组成\term[jiben-qun]{基本群}[fundamental group] \(\pi_1(U,a)\)，乘法由路径拼接给出。
\end{definition}
前一命题保证这确为群。区域路径连通，换基点的连接路径给群同构，但同构依赖所选连接路径，两个选择可能相差一个内自同构。单连通等价于任意基点的基本群平凡。

\ProseParagraphBreak
凸区域中的闭路由直线收缩 \(H(s,t)=(1-s)\gamma(t)+sa\)，故基本群平凡。若区域仅星形，同样使用星形中心，再用基点连接路径换基点。开集连通保证路径连通，可以用“能由有限内部线段到达的点集同时开且闭”证明，不需要假定区域本身凸。

\ProseParagraphBreak
换基点时，群同构并非完全没有选择。若 \(\eta\) 从 \(a\) 到 \(b\)，把基于 \(b\) 的闭路 \(\gamma\) 送到 \(\eta*\gamma*\eta^{-1}\)，得到基于 \(a\) 的类。换另一条 \(\eta\)，所得同构相差一个内自同构。因而只讨论群的抽象同构类型时可不记基点，讨论具体闭路对芽的作用时却须记录基点及连接路径。两处都写“绕某个孔一圈”，若未说明如何从基点走到孔附近，可能代表不同的共轭类。

\ProseParagraphBreak
拼接路径的参数变化不会影响其同伦类。把每段路径分配不同的时间长度，线性插值各分段时刻即可构造同伦；插入常路径也可用同样办法消去。逆路径使 \(\gamma*\gamma^{-1}\) 可收缩，具体同伦让去程只走到 \(\gamma(1-s)\) 后立即沿原路返回。群公理因此来自可写出的路径变形，而不是依赖参数拼接严格满足结合律。

\subsection{绕行}
\label{b3:subsec:090104}

\begin{proposition}\label{b3:prop:09-punctured-fundamental}
\(\pi_1(\mathbb C^*,1)\) 由绕数同构于 \(\mathbb Z\)。
\end{proposition}
\begin{proof}
径向变形 \(z\mapsto ((1-s)|z|+s)z/|z|\) 将去心平面收缩到单位圆并固定单位圆。圆周路径有连续辐角提升 \(\gamma(t)=e^{\mathrm{i}\theta(t)}\)，可沿有限小弧依次选择，指定 \(\theta(0)=0\) 后唯一。闭路给 \(\theta(1)=2\pi n\)。端点固定同伦中的终点整数连续不变；反之具有相同 \(n\) 的两个辐角提升可作线性插值，指数后给端点固定同伦。因此类由整数唯一决定，拼接使整数相加。
\end{proof}
绕两个不同点的区域不再由一个整数分类。绕数只记录每个点的总绕行次数，而不记录绕行顺序；两个孔的基本群中交换子可以非平凡，却对每个外点绕数为零。第10章将区分这类同伦信息与较弱的同调信息。
